\documentclass[a4paper, 10pt]{article}
\usepackage{scrextend}
\changefontsizes{9pt}
% Per-subject config for this repo. See vendor/template/course-config.tex.example
% for what each macro is used for.

\newcommand{\CourseCode}{2110506}
\newcommand{\CourseName}{Software-Defined Systems}
\newcommand{\StudentID}{6733172621}
\newcommand{\StudentName}{Patthadon Phengpinij}
\newcommand{\DefaultCollaborators}{Claude (for\,\LaTeX\,styling and grammar checking)}
\newcommand{\AssignmentLabel}{Activity}

\usepackage{../vendor/template/CEDT-Assignment-style}

\begin{document}
\subject
\asgntitle{4}{}{Week 5}{\StudentID\ \StudentName}{\DefaultCollaborators}

% ================================================================================ %
\section*{Activity: Kubernetes with Multiple Containers}
% ================================================================================ %

\textbf{Instructions}
\begin{enumerate}
	\item This assignment is due Wednesday September 9th no later than 11:00 pm on mycourseville (pdf only).
	\item Install a kubernetes cluster (using docker desktop or k3d or other distributions).
	\item Test your cluster with nginx as followed:
	      \begin{subproblems_alpha}
		      \item Deploy an nginx on your cluster using the deployment file from
		            \par\href{https://k8s.io/examples/application/deployment.yaml}{https://k8s.io/examples/application/deployment.yaml}
		      \item Use kubectl to describe the deployment
	      \end{subproblems_alpha}
	\item Deploy a \texttt{todo} service with redis database as followed:
	      \begin{subproblems_alpha}
		      \item Study how to create a deployment of multiple containers in a single pod from
		            \par\small{\href{https://www.mirantis.com/blog/multi-container-pods-and-container-communication-in-kubernetes/}{https://www.mirantis.com/blog/multi-container-pods-and-container-communication-in-kubernetes/}}
		      \item Create a deployment file of a pod with service of the following containers:
		            \begin{enumerate}[label=\roman*)]
			            \item \texttt{todo} (release-2.1)
			            \item \texttt{redis}
		            \end{enumerate}
		            \par \textit{Note:} you can use a docker-hub image: \texttt{natawut/todo-service:release-2.1} and you must ensure that environment for \texttt{REDIS\_HOST} is correctly set.
		      \item Deploy the pod, expose service, and show your deployment, service, and pods with kubectl get command.
		      \item Create ingress configuration file, apply the file, and show the ingress information.
		      \item Test the newly deployment:
		            \begin{enumerate}[label=\roman*)]
			            \item Create a new todo record with:
			                  \begin{codingbox}[language=bash]
{"title":"todo-1","detail":"the first todo","duedate":"2022-10-24 11:00:00","tags":[],"completed":false}
\end{codingbox}
			            \item Get \texttt{todo} list via \href{http://localhost:80}{localhost:80}
		            \end{enumerate}
		            \par Below shows the sample output:
		            \begin{codingbox}[language=bash]
curl -X POST -d '{"title":"todo-1","detail":"the first todo","duedate":"2024-10-13 11:00:00","tags":[],"completed":false}' http://localhost
\end{codingbox}
		            \fig[\linewidth]{images/step-04-example.png}{Example output of the \texttt{todo} service.}{step-04-example}
		            \par \textit{Note:} to help you debug, you may use kubectl logs to check the system log in the container
	      \end{subproblems_alpha}
	\item Use the activity grader here:
	      \par\href{https://github.com/fonylew/sds-grader/releases#release-v0.3.3-activity4}{https://github.com/fonylew/sds-grader/releases\#release-v0.3.3-activity4}
	\item Prepare a PDF report containing the following information:
	      \begin{subproblems_alpha}
		      \item Screenshot when you perform Step 3.
		      \item The deployment file that you create in Step 4.b and ingress file in Step 4.d.
		      \item Screenshot when you perform Step 4.c - Step 4.e.
	      \end{subproblems_alpha}
\end{enumerate}

\begin{center}
	\par\noindent\textbf{\color{red} YOU MUST ALSO USING GRADER (URL is in MCV's assignment)}
\end{center}

\newpage

\noindent\textbf{Environment.}
The cluster is a local single-node \textbf{k3d} cluster (k3s in Docker), created with
host port \texttt{80} mapped to the cluster load balancer so the ingress is reachable
at \texttt{http://localhost}; k3d ships \textbf{Traefik} as the default ingress
controller. \texttt{kubectl}~v1.37 against server~v1.35.
\begin{codingbox*}[language=bash]
$ k3d cluster create sds --port "80:80@loadbalancer"
$ kubectl get nodes
NAME               STATUS   ROLES                  AGE   VERSION
k3d-sds-server-0   Ready    control-plane,master   1m    v1.35.5+k3s1
\end{codingbox*}

\bigskip

% ================================================================================ %
%                                    Problem 01                                    %
% ================================================================================ %
\begin{problem}
Screenshot when you perform Step 3.
\end{problem}

\begin{solution}
	The example deployment is applied with
	\texttt{kubectl apply -f https://k8s.io/examples/application/deployment.yaml}
	and then inspected with \texttt{kubectl describe deployment nginx-deployment}.
	It runs 2 replicas of \texttt{nginx:1.14.2}, both \texttt{Available}.

	\fig[\linewidth]{images/step-03.png}{Step 3 --- deploy the example nginx and \texttt{kubectl describe deployment nginx-deployment}.}{step-03}
\end{solution}
% ================================================================================ %

% ================================================================================ %
%                                    Problem 02                                    %
% ================================================================================ %
\begin{problem}
The deployment file that you create in Step 4.b and ingress file in Step 4.d.
\end{problem}

\begin{solution}
	\textbf{(a) Deployment --- \texttt{todo-deployment.yaml} (Step 4.b).}
	One pod with two containers, \texttt{todo} and \texttt{redis}. Containers in the
	same pod share a network namespace, so the todo service reaches redis on
	\texttt{localhost} --- hence \texttt{REDIS\_HOST: localhost}.
	\begin{codingbox*}[language=YAML]
apiVersion: apps/v1
kind: Deployment
metadata:
  name: todo
  labels:
    app: todo
spec:
  replicas: 1
  selector:
    matchLabels:
      app: todo
  template:
    metadata:
      labels:
        app: todo
    spec:
      containers:
        - name: todo
          image: natawut/todo-service:release-2.1
          imagePullPolicy: IfNotPresent
          ports:
            - containerPort: 8000
          env:
            - name: REDIS_HOST
              value: "localhost"      # same pod -> same network namespace
            - name: REDIS_PORT
              value: "6379"
        - name: redis
          image: redis:7-alpine
          imagePullPolicy: IfNotPresent
          ports:
            - containerPort: 6379
	\end{codingbox*}

	\textbf{(b) Service --- \texttt{todo-service.yaml} (Step 4.c, ``expose service'').}
	A ClusterIP service in front of the todo container's port \texttt{8000}. It is
	\emph{not} exposed on the host, so \texttt{localhost:8000} and \texttt{localhost:6379}
	stay unreachable from outside --- only the ingress is public.
	\begin{codingbox*}[language=YAML]
apiVersion: v1
kind: Service
metadata:
  name: todo
  labels:
    app: todo
spec:
  selector:
    app: todo
  ports:
    - name: http
      port: 8000
      targetPort: 8000
	\end{codingbox*}

	\textbf{(c) Ingress --- \texttt{todo-ingress.yaml} (Step 4.d).}
	Routes \texttt{http://localhost/} to the todo service. The resource is named
	\texttt{todo-ingress} (the grader runs \texttt{kubectl get ingress} and only
	accepts the result when the word ``ingress'' appears in that output).
	\begin{codingbox*}[language=YAML]
apiVersion: networking.k8s.io/v1
kind: Ingress
metadata:
  name: todo-ingress
spec:
  ingressClassName: traefik
  rules:
    - host: localhost
      http:
        paths:
          - path: /
            pathType: Prefix
            backend:
              service:
                name: todo
                port:
                  number: 8000
	\end{codingbox*}
\end{solution}
% ================================================================================ %

% ================================================================================ %
%                                    Problem 03                                    %
% ================================================================================ %
\begin{problem}
Screenshot when you perform Step 4.c - Step 4.e.
\end{problem}

\begin{solution}
	Step 4.c applies \texttt{todo-deployment.yaml} and \texttt{todo-service.yaml} and shows
	\texttt{kubectl get deployment,svc,pods} --- the pod is \texttt{2/2} ready (both the
	\texttt{todo} and \texttt{redis} containers). Step 4.d applies \texttt{todo-ingress.yaml}
	and shows \texttt{kubectl get/describe ingress}. Step 4.e \texttt{POST}s a todo to
	\texttt{http://localhost} (returns \texttt{\{"id":0\}}) and \texttt{GET}s the list back.

	\fig[\linewidth]{images/step-04.png}{Step 4.c--4.e --- \texttt{kubectl get deployment,svc,pods}, \texttt{describe ingress}, and the \texttt{curl} round-trip against \texttt{http://localhost}.}{step-04}
\end{solution}
% ================================================================================ %

% ================================================================================ %
%                                    Problem 04                                    %
% ================================================================================ %
\begin{problem}
Screenshot of the results after running the \texttt{GRADER}.
\end{problem}

\begin{solution}
	The \texttt{activity4-CEDT} grader passes every check --- \texttt{Result: true} ---
	then prompts for the Student ID and full name to submit the result.

	\fig[\linewidth]{images/grader.png}{\texttt{activity4-CEDT} grader --- all checks pass, \texttt{Result: true}.}{grader}
\end{solution}
% ================================================================================ %

\end{document}
